% !TeX root = ../../../../../main.tex
% chapters/sections/chapter5/subsections/asad/ad.tex

\subsection{政策による均衡の制御と期待の役割}
\label{sec:chapter5_asad_monetary_policy}
本モデルの均衡は水平なAS曲線
\eqref{form:chapter5_asad_as_p_H_eq}
と右下がりのAD曲線
\eqref{form:chapter5_asad_ad_p_H_eq}
の交点として描かれる。
先に述べたように\(\mathcal{B}_t\)と\(\Lambda_{\infty}\)は金融政策の影響を受けないため、
金融政策の分析上考慮すべき変数は
\(p_t^H,y_t^H,i_t^H,E_t[\lambda_{t+1}^H],E_t[p_{t+1}^H]\)
の5つである。
以下では\(i_t^H,E_t[\lambda_{t+1}^H],E_t[p_{t+1}^H]\)を仮想的に単独で変化させ
\((y,p)\)平面上の均衡点の動きを観察する。

\subsubsection{名目利子率\(i_t^H\)の低下}
\label{sec:chapter5_asad_monetary_policy_i_H_fall}
名目利子率\(i_t^H\)が低下するとAD曲線
\eqref{form:chapter5_asad_ad_p_H_eq}
の右辺分母が小さくなり右辺全体が大きくなるため、
\textbf{AD曲線は右上へ移動}する。
これにより生産\(y_t^H\)は増加する。

\subsubsection{所得の期待限界効用\(E_t[\lambda_{t+1}^H]\)の低下}
\label{sec:chapter5_asad_monetary_policy_E_lambda_H_fall}
所得の期待限界効用\(E_t[\lambda_{t+1}^H]\)が低下すると
AD曲線
\eqref{form:chapter5_asad_ad_p_H_eq}
の右辺分母が小さくなり右辺全体が大きくなるため、\textbf{AD曲線は右上へ移動}する。
これにより生産\(y_t^H\)は増加する。
これは名目利子率\(i_t^H\)を下げる余地がないゼロ金利制約下にあっても、
中央銀行の約束が家計の期待に働きかけることにより
生産\(y_t^H\)を回復させることが可能であることを示している。

\subsubsection{期待生産者物価指数\(E_t[p_{t+1}^H]\)の上昇}
\label{sec:chapter5_asad_monetary_policy_E_p_H_rise}
期待生産者物価指数\(E_t[p_{t+1}^H]\)が上昇した場合、
AS曲線
\eqref{form:chapter5_asad_as_p_H_eq}
の右辺第2項が増大するため、\textbf{AS曲線は上へ移動}する。
これにより生産者物価指数\(p_t^H\)は上昇し、生産\(y_t^H\)は下落する。
しかし、このとき期待限界効用\(E_t[\lambda_{t+1}^H]\)は
所得の限界効用の式
\eqref{form:chapter3_representative_household_lambda_home_representative_lambda_H_p_H_W_c_H_W_eq}
を通じて減少するため（TODO）、先に述べたように\textbf{AD曲線は右上へ移動}し、
これは生産\(y_t^H\)を増加させる方向に働く。

\subsubsection{名目総消費水準目標の働きかけ}
名目総消費水準目標はこれら3つの指標に総合的に働きかけることにより、
ゼロ金利制約に直面した経済を効率的に回復させる。